\secciongrande{Herramientas de programación sobre clientes web: Elección del editor de código}

\resumen{
  La elección del editor de código depende de las necesidades del desarrollador, priorizando herramientas de desarrollo sobre las de modificación de contenido. Es crucial que ofrezcan funcionalidades como codificación de colores, verificación de sintaxis y soporte FTP. Editores gratuitos recomendados incluyen Notepad++, Eclipse, NetBeans, KompoZer y Bluefish.
}

\definicion{Editor de código}{Herramienta especializada para la escritura y edición de código fuente, optimizada para lenguajes de programación como HTML y JavaScript, a diferencia de los editores orientados a la presentación.}
\subtituloA{Criterios de selección de herramientas}

\lista{
  \textbf{Exclusión:} No utilizar editores orientados a modificación de contenidos y presentación (ej. FrontPage, Dreamweaver). |
  \textbf{Prioridad:} Facilidad para almacenar archivos HTML y desarrollo web, incluyendo codificación de colores, verificación de sintaxis, inserción de comentarios y soporte FTP.
}
\nota{Para el examen}{
Diferenciar entre un desarrollador de código y una herramienta de modificación de contenidos es fundamental para el desarrollo web profesional.
}
